\section{Implementation}
\subsection{Proposed Board Connection}
The example target is Arduino UNO R3. Its digital inputs and outputs are enough
for this small control unit~\cite{uno}. Three extra switches emulate coin signals;
they are test inputs, not extra machine controls. An optional fault switch
provides the error signal.

\begin{tabularx}{\linewidth}{@{}p{3.5cm}p{2.05cm}Y@{}}
\toprule
Signal & Pin & Connection\\
\midrule
STOP, RUN, PAUSE & D2, D3, D4 & Each switch connects its input to GND.\\
10, 20, 50 cent test inputs & D5, D6, D7 & Each switch models one valid coin insertion.\\
RLED & D8 & D8 to 330 $\Omega$ resistor, LED anode, then cathode to GND.\\
BLED & D9 & D9 to 330 $\Omega$ resistor, LED anode, then cathode to GND.\\
Wash enable & D10 & Digital enable for a separate driver; high only in Running.\\
Fault test input & A0 & Switch to GND; LOW enters Error.\\
Power and reference & USB, GND & USB supplies the board; all test inputs share GND.\\
\bottomrule
\end{tabularx}

\fig{03_wiring_diagram.pdf}{1.0}{Proposed signal connections for the bench circuit.}

All switch inputs use \code{INPUT\_PULLUP}. An open switch reads HIGH and a
pressed switch reads LOW~\cite{pullup}. Each button must stay stable for 30 ms
before a press is counted. Release is also debounced, so a held STOP cannot be
mistaken for two presses.

The bench parts are one UNO R3, two LEDs, two 330 $\Omega$ resistors, six push
buttons, one optional fault switch, a breadboard, jumper wires and a USB cable.
Coins can also be entered through Serial Monitor, so the three coin switches
are optional for a simple demonstration.

D10 represents the control decision. A motor requires a separate driver and
power supply. No motor or mains circuit is needed to show the required state logic.
